%%%%%%%% ICML 2026 LATEX FILE — Solve Vision with Code %%%%%%%%%%%%%%%%%
% Adapted from the object-permanence paper template; preamble kept, content replaced.
% Build: bash build.sh   (pdflatex → bibtex → pdflatex → pdflatex)

\documentclass{article}
\usepackage{hyphenat}
% Recommended, but optional, packages for figures and better typesetting:
\usepackage{microtype}
\usepackage{graphicx}
\graphicspath{{figures/}{figures/qualitative/}{sec/}} % figures/*.pdf, qualitative/*, legacy sec/* all resolve
\usepackage{subcaption}
\usepackage{booktabs} % for professional tables
\usepackage{multirow}
\usepackage{longtable}
\usepackage{float}
\usepackage{caption}

% TIGHT UP SPACE
 %%%%%%%%%%%%%
% \setlength{\parskip}{0pt}
% \setlength{\parindent}{1em}  % keep a
% \usepackage{titlesec}
% \titlespacing*{\paragraph}{0pt}{0.3ex plus 0.1ex minus 0.1ex}{0.6em} %normal indent (optional)
% \makeatletter
% \renewcommand\paragraph{\@startsection{paragraph}{4}{\z@}%
%  {0.35ex \@plus 0.1ex \@minus 0.1ex}% space above
 % {-0.8em}% space after (negative = run-in)
 % {\normalfont\normalsize\bfseries}}
% \makeatother

% hyperref makes hyperlinks in the resulting PDF.
% If your build breaks (sometimes temporarily if a hyperlink spans a page)
% please comment out the following usepackage line and replace
% \usepackage{icml2026} with \usepackage[nohyperref]{icml2026} above.
\usepackage{hyperref}
\usepackage[ruled,vlined]{algorithm2e}

% Attempt to make hyperref and algorithmic work together better:
\usepackage{wrapfig}\usepackage[table]{xcolor} % [table] already loads colortbl
\definecolor{bestcolor}{RGB}{219, 208, 237}
\definecolor{secondcolor}{RGB}{241, 237, 248}
\definecolor{line-blue}{RGB}{243, 248, 252}
\definecolor{taskSpatial}{RGB}{255, 204, 170}   % Peach
\definecolor{taskTransform}{RGB}{210, 190, 230} % Lavender
\definecolor{taskLogic}{RGB}{160, 200, 255}     % Sky Blue
\definecolor{taskAbstract}{RGB}{245, 170, 185}  % Rose Pink
\definecolor{taskPercept}{RGB}{180, 225, 200}   % Soft Mint
\definecolor{taskPhysics}{RGB}{30, 60, 120}

\definecolor{resourceblue}{HTML}{4C7DBF}
\newcommand{\pending}[1]{\textcolor{red}{[#1]}}
\newcommand{\resource}[2]{\href{#2}{\textcolor{resourceblue}{[\textbf{#1}]}}}

\newcommand{\theHalgorithm}{\arabic{algorithm}}

% Macros
\usepackage{xspace}



% Use the following line for the initial blind version submitted for review:
% \usepackage{icml2026}

% For preprint, use
% \usepackage[preprint]{icml2026}

% THIS IS FOR PUBLIC RELATIONS
\usepackage[pr]{icml2026}

\makeatletter
% Header brand mark: icml2026.sty [pr] hard-codes the object-permanence site; point it at this repo.
\renewcommand{\pr@headerlogo}{%
  {\pr@logofont\textcolor{pr@accent}{%
    \href{https://hokindeng.com}{hokindeng.com}}}%
}
\gdef\affil@sep{\enspace}
\newcommand{\affilbreak}{\gdef\affil@sep{\\[2pt]}}

\renewcommand{\printAffiliationsAndNotice}[1]{\global\icml@noticeprintedtrue%
  \stepcounter{@affiliationcounter}%
  \noindent\begin{minipage}{\textwidth}%
    \centering
    \fontsize{9}{11}\selectfont%
    \forloop{@affilnum}{1}{\value{@affilnum} < \value{@affiliationcounter}}{%
      \textsuperscript{\arabic{@affilnum}}\,\csname @affilname\the@affilnum\endcsname%
      \affil@sep\gdef\affil@sep{\enspace}%
    }%
    \ificmlshowauthors #1\fi%
    \ifdefined\icmlcorrespondingauthor@text
      \\[4pt]%
      \textsuperscript{\textdagger}\,Correspondence to: \icmlcorrespondingauthor@text.%
    \fi%
    \\[6pt]% resource links start here
    \resource{Code}{https://github.com/hokindeng/solve-vision-with-code}\enspace
    \resource{Website}{https://hokindeng.com}\enspace
    \resource{Results}{https://hokindeng.github.io/solve-vision-with-code/}%
  \end{minipage}%
  \vskip 8pt%
  {\centering\textcolor{pr@lightgray}{\rule{0.2\textwidth}{0.3pt}}\par}%
}
\makeatother

% If accepted, instead use the following line for the camera-ready submission:
% \usepackage[accepted]{icml2026}

\usepackage{amsmath}
\usepackage{amssymb}
\usepackage{mathtools}
\usepackage{amsthm}
% xcolor already loaded above with [table] option

% if you use cleveref..
\usepackage[capitalize,noabbrev]{cleveref} % LAST

%%%%%%%%%%%%%%%%%%%%%%%%%%%%%%%%
% THEOREMS
%%%%%%%%%%%%%%%%%%%%%%%%%%%%%%%%
\theoremstyle{plain}
\newtheorem{theorem}{Theorem}[section]
\newtheorem{proposition}[theorem]{Proposition}
\newtheorem{lemma}[theorem]{Lemma}
\newtheorem{corollary}[theorem]{Corollary}
\theoremstyle{definition}
\newtheorem{definition}[theorem]{Definition}
\newtheorem{assumption}[theorem]{Assumption}
\theoremstyle{remark}
\newtheorem{remark}[theorem]{Remark}

% Todonotes is useful during development; simply uncomment the next line
%    and comment out the line below the next line to turn off comments
\usepackage[disable,textsize=tiny]{todonotes}
%\usepackage[textsize=tiny]{todonotes}
\usepackage{enumitem}

\usepackage{makecell}
\usepackage{verbatim}   % \verbatiminput of solve.py listings in the appendix / analysis

% Float placement: the single-column [pr] layout otherwise defers most tables/figures to after the
% conclusion. Let a page be mostly float and allow more floats per page.
\renewcommand{\topfraction}{0.95}
\renewcommand{\bottomfraction}{0.7}
\renewcommand{\textfraction}{0.05}
\renewcommand{\floatpagefraction}{0.8}
\setcounter{topnumber}{3}
\setcounter{bottomnumber}{2}
\setcounter{totalnumber}{5}

% Short form of the title for the running header:
\icmltitlerunning{Solve Vision with Code}

\begin{document}
\twocolumn[
  \icmltitle{Solve Vision with Code: Coding Agents Are Stronger Visual Reasoners Than Video Models}

  % It is OKAY to include author information, even for blind submissions: the
  % style file will automatically remove it for you unless you've provided
  % the [accepted] option to the icml2026 package.

  % List of affiliations: The first argument should be a (short) identifier you
  % will use later to specify author affiliations Academic affiliations
  % should list Department, University, City, Region, Country Industry
  % affiliations should list Company, City, Region, Country

  % You can specify symbols, otherwise they are numbered in order. Ideally, you
  % should not use this facility. Affiliations will be numbered in order of
  % appearance and this is the preferred way.
  % \icmlsetsymbol{equal}{*}


\begin{icmlauthorlist}
    \icmlauthor{Hokin Deng}{cmu}
    \icmlauthor{Zehong Zhao}{ucsd}
    \icmlauthor{Ran Ji}{ucsd}
    \icmlauthor{Maijunxian Wang}{ucb}
    \icmlauthor{Qingying Gao}{jhu}
    \icmlauthor{Fengyuan Yu}{cmu}
    \icmlauthor{Zhengze Jiang}{columbia}
    \icmlauthor{Yilan Zhang}{ucla}
\end{icmlauthorlist}

\icmlaffiliation{cmu}{Carnegie Mellon University}
\icmlaffiliation{ucsd}{University of California, San Diego}
\icmlaffiliation{ucb}{University of California, Berkeley}
\icmlaffiliation{jhu}{Johns Hopkins University}
\icmlaffiliation{columbia}{Columbia University}
\icmlaffiliation{ucla}{University of California, Los Angeles}

\icmlcorrespondingauthor{Hokin Deng}{hokind@andrew.cmu.edu}
\printAffiliationsAndNotice{}
% \printAffiliationsAndNotice{\icmlEqualContribution}

\icmlkeywords{Visual Reasoning, Coding Agents, Video Models, Program Synthesis}
]

% this must go after the closing bracket ] following \twocolumn[ ...

% This command actually creates the footnote in the first column listing the
% affiliations and the copyright notice. The command takes one argument, which
% is text to display at the start of the footnote. The \icmlEqualContribution
% command is standard text for equal contribution. Remove it (just {}) if you
% do not need this facility.

% Use ONE of the following lines. DO NOT remove the command.
% If you have no special notice, KEEP empty braces:
% \printAffiliationsAndNotice{}  % no special notice (required even if empty)
% Or, if applicable, use the standard equal contribution text:

% Each section is pulled in only if its file exists; a missing one leaves a red note
% instead of breaking the build (the section files are written by separate owners).
\newcommand{\inputsection}[2]{\IfFileExists{#1.tex}{\input{#1}}{\section*{\pending{section #2 missing: \texttt{\detokenize{#1.tex}}}}}}

\inputsection{sec/0_abstract}{0 (abstract)}
\inputsection{sec/1_introduction}{1 (introduction)}
\inputsection{sec/2_related_works}{2 (related work)}
\inputsection{sec/3_method}{3 (method)}
\inputsection{sec/4_results}{4 (results)}
\inputsection{sec/5_analysis}{5 (analysis)}
\inputsection{sec/6_conclusion}{6 (conclusion)}

\clearpage

% In the unusual situation where you want a paper to appear in the
% references without citing it in the main text, use \nocite
% \nocite{langley00}

\IfFileExists{main.bib}{%
  \bibliography{main}
  \bibliographystyle{icml2026}
}{\section*{\pending{main.bib missing: references not built}}}

\clearpage
\appendix
\inputsection{sec/appendix}{A (appendix)}

\end{document}
